\documentclass[11pt]{article}
\usepackage{amsmath,amssymb,amsthm}
\usepackage[margin=1in]{geometry}
\usepackage{hyperref}
\title{Disproof of TLMC Conjecture 00000001068}
\author{}
\date{October 2, 2026}
\newtheorem{theorem}{Theorem}
\newtheorem{lemma}[theorem]{Lemma}
\newcommand{\F}{\mathbb{F}}
\begin{document}
\maketitle

\section*{Verdict}
\textbf{FALSE.} The conjecture fails already at $p = 11$.

\section{The conjecture}

\begin{quote}
\textit{The maximal sum-free subsets of $\F_p$ are exactly the intervals
$((p+1)/3,\, 2(p-1)/3)$, uniquely up to dilation (the complete
Diananda--Yap characterization over finite fields).}
\end{quote}

A set $A \subseteq \F_p$ is \emph{sum-free} if $(A + A) \cap A = \emptyset$,
and \emph{maximal} sum-free if it is not properly contained in another
sum-free set.

\section{Counterexample: \texorpdfstring{$p = 11$}{p = 11}}

\begin{theorem}\label{thm:main}
In $\F_{11}$ the set $A = \{4,5,6,7\}$ is a maximal sum-free subset with
$|A| = 4$, whereas the conjectured interval
$I = ((p+1)/3,\, 2(p-1)/3) = (4,\, 20/3) = \{5,6\}$ has $|I| = 2$ and is not
maximal. Hence the conjecture is false.
\end{theorem}

\begin{proof}
\emph{$A$ is sum-free.} A direct computation of all sums of two elements of
$A$ gives
\[
A + A \;=\; \{0,1,2,3,8,9,10\} \;=\; \F_{11} \setminus A,
\]
which is disjoint from $A$.

\emph{$A$ is maximal.} Every element outside $A$ is a sum of two elements of
$A$:
\[
0 = 5+6,\quad 1 = 5+7,\quad 2 = 6+7,\quad 3 = 7+7,\quad
8 = 4+4,\quad 9 = 4+5,\quad 10 = 4+6 .
\]
If $x \notin A$ were added to $A$, then $x \in A + A \subseteq
(A \cup \{x\}) + (A \cup \{x\})$ would meet $A \cup \{x\}$, so no enlargement
of $A$ is sum-free.

\emph{$I$ is not maximal.} By definition of the open interval with exact
rational endpoints, $x \in I$ iff $3x > 12$ and $3x < 20$, i.e.\
$I = \{5,6\}$. Since $I \subsetneq A$ and $A$ is sum-free, $I$ is not maximal.
The same holds for every nonzero dilate: for $d = 1,\dots,10$ the dilate
$dI = \{5d, 6d\} \bmod 11$ has exactly two elements and satisfies
$dI \subsetneq dA$, where $dA$ is sum-free (dilation by a unit preserves
sum-freeness, and this is directly verified for all $d$); the element
$4d \bmod 11$ lies in $dA \setminus dI$. Thus the family conjectured to be
``the maximal sum-free subsets'' contains no maximal sum-free set at
$p = 11$.

\emph{$A$ lies outside the conjectured family.} $A$ is maximal sum-free by
the above, and every dilate of $I$ has $2$ elements while $|A| = 4$; so $A$
is not a dilate of $I$.
\end{proof}

Both inclusions of the claimed characterization therefore fail: the
conjectured family contains no maximal sum-free set, and there exists a
maximal sum-free set outside it.

\section{Computational verification}

A brute-force census over all $2^{11}$ subsets of $\F_{11}$ confirms the
maximum sum-free size is $4 = \lfloor (p+1)/3 \rfloor$ and that there are
exactly $15$ maximal sum-free subsets, with sizes in $\{3,4\}$;
$A = \{4,5,6,7\}$ is one of the size-$4$ maximal sets. All numbers above are
recomputed by the accompanying script \texttt{reproduce.py} and are
machine-checked in pure Lean~4 (no axioms, no \texttt{sorry}; see
\texttt{lean4/}).

\section{Boundary remarks}

\begin{itemize}
\item The attack is robust to endpoint conventions: with integer-division
endpoints the open reading gives $(4,6) = \{5\}$ and the closed reading gives
$[4,6] = \{4,5,6\}$; both are still strictly contained in the sum-free set
$A$, hence non-maximal under every reading.
\item The breakage is not specific to $p = 11$. Under the exact-rational open
reading, $|I|$ equals $0,0,0,1,2,3$ for $p = 2,3,5,7,11,13$, while the true
maximum sum-free size equals $1,1,2,2,4,4$; the claimed interval is
non-maximal for every prime in this census.
\item The genuine Diananda--Yap extremal interval is
$\{\lfloor p/3 \rfloor + 1, \dots, \lfloor 2p/3 \rfloor\}$, equal to
$\{4,\dots,7\}$ for $p = 11$. The conjecture's open endpoints
$((p+1)/3,\, 2(p-1)/3)$ truncate that interval, which is precisely what
destroys maximality. A single counterexample refutes the universal claim.
\end{itemize}

\end{document}
